% Minimal wrapper for the Saar–Koyama §X technical/computational
% section. Designed for drop-in compilation while the host paper
% structure is being assembled; on integration the host paper will
% replace this preamble and \input the three subfiles.
%
% Compile: pdflatex paper && bibtex paper && pdflatex paper && pdflatex paper

\documentclass[11pt]{amsart}

% ---------- packages ---------------------------------------------------------
\usepackage[margin=1in]{geometry}
\usepackage[utf8]{inputenc}
\usepackage[T1]{fontenc}
\usepackage{lmodern}
\usepackage{microtype}

\usepackage{amsmath, amssymb, amsthm, mathtools}
\usepackage{booktabs, longtable, array}
\usepackage{enumitem}
\usepackage{etoolbox}
% Wide tables (esp. §X.6 Lean inventory, §X.5 numerical residuals)
% overflow the text width at \normalsize. Render them at \small with
% tightened column spacing and a generous \emergencystretch so TeX
% can find acceptable line breaks inside narrow cells.
\AtBeginEnvironment{longtable}{\small\setlength{\emergencystretch}{3em}}
\setlength{\tabcolsep}{4pt}
\renewcommand{\arraystretch}{1.05}
\usepackage[dvipsnames]{xcolor}
\usepackage{hyperref}
\hypersetup{colorlinks=true, linkcolor=MidnightBlue, citecolor=MidnightBlue, urlcolor=MidnightBlue}

% ---------- theorem environments --------------------------------------------
\newtheorem{theorem}{Theorem}[section]
\newtheorem{lemma}[theorem]{Lemma}
\newtheorem{proposition}[theorem]{Proposition}
\newtheorem{corollary}[theorem]{Corollary}

\theoremstyle{definition}
\newtheorem{definition}[theorem]{Definition}
\newtheorem{conjecture}[theorem]{Conjecture}
\newtheorem{question}[theorem]{Question}
\newtheorem{remark}[theorem]{Remark}
\newtheorem*{hypothesis}{Hypothesis}

% ---------- macros used in §X / appendices ----------------------------------
% Note: the markdown source writes \mathrm{Re}(...) directly, so we don't
% redefine \Re.
\newcommand{\Res}{\operatorname{Res}}
\newcommand{\BPC}{\mathrm{BPC}}
\newcommand{\dps}{\mathrm{dps}}

% ---------- bibliography ----------------------------------------------------
\usepackage[backend=bibtex,style=alphabetic,maxbibnames=4]{biblatex}
\addbibresource{references.bib}

% ---------- meta -------------------------------------------------------------
\title[Technical / Computational section]{Methodology, formalization, and numerical evidence \\
       \large (Technical / Computational section --- draft)}
\author{Shin-ya Koyama \and Saar Shai}
\date{\today}

% Allow TeX more flexibility with line breaks. Long inline-code
% tokens (e.g. `formal-conjectures/LocalPerronResidue.lean`) and code
% identifiers are otherwise unbreakable and cause overfull boxes.
\sloppy
\setlength{\emergencystretch}{4em}
\hyphenpenalty=1000
\exhyphenpenalty=1000
% Presentation-only line-breaking tolerance (does not alter content):
% lets TeX accept slightly looser lines instead of emitting overfull
% \hbox rules inside the dense \S X.5/\S X.6 tables.
\tolerance=2000
\hbadness=2000
% Permit breaks inside long unbreakable code tokens (Lean module
% paths/identifiers) as a last resort, which is the dominant source
% of the residual overfull boxes in the inventory tables.
\appto\ttfamily{\hyphenchar\font=`\-}

\begin{document}

\begin{abstract}
This is the technical and computational section (\S X) and its two
proof appendices of the joint manuscript, circulated as a
self-contained draft. It records the corrected $B_\infty$
framework: the four-component $B_\infty$ identity
(Theorem~X.4.1, unconditional given simplicity of the zero
$\rho$; full proof in Appendix~A), the leading-plus-%
subleading partial-M\"obius identity (Theorem~X.4.2, Appendix~B),
the local Perron double-pole residue (Lemma~X.3.1), a double-%
verified numerical study at two rigorously separated scales (the
Phase-1 Dominance-of-$-1$ replication at $x = 1.3\cdot10^{13}$ and
the analytic-identity verification to $K = 10^{8}$), and a Lean~4
/ Mathlib formalisation inventory with a cumulative axiom audit.
The full paper's abstract and introduction are drafted jointly and
are not part of this technical draft.
\end{abstract}

\maketitle

\section{Methodology, formalization, and numerical evidence}

\emph{The section number \texttt{X} and the equation / theorem / lemma tags
(Lemma X.3.1, Theorem X.4.1, etc.) are placeholders to be assigned
on integration into the full paper.}

This section records the technical and computational content of the
paper: the algebraic identities at the heart of the corrected
\(B_\infty\) framework (\S\,X.3--\S\,X.4), the open challenges that organise
the program forward (\S\,X.4.4, \S\,X.7), the numerical evidence at the
two scales on which our verification operates (\S\,X.5), and the Lean 4
formalisation inventory (\S\,X.6).

The numerical evidence has \textbf{two distinct verified scales}, which
we keep rigorously separate throughout. The \textbf{replication scale}
is \(x = 1.3 \cdot 10^{13}\), at which two independent prime-counting
implementations agree on Koyama's Dominance-of-\(-1\) residue tables
(\S\,X.5.1). The \textbf{analytic-identity scale} is \(K \le 10^{8}\), at
which the corrected \(B_\infty\) identity (\S\,X.5.4), the subleading
constant \(C_1\) (\S\,X.5.2), and the Aoki--Koyama--Mertens drift toward
\(e^{-\gamma}\) (\S\,X.5.3, at \(K \le 10^{7}\)) are verified across
three numerical stacks. Numbers
from one scale are not transferred to the other.

\subsection{Notation}

Fix a primitive non-principal Dirichlet character \(\chi\) modulo
\(q \ge 2\), and let \(\rho = \tfrac12 + i\tau\) with \(\tau \neq 0\) be a
simple non-trivial zero of \(L(s,\chi)\). We use:

\begin{itemize}
\item
  \(E_K(\chi,\rho) := \prod_{p \le K} (1 - \chi(p)\,p^{-\rho})^{-1}\)
  (truncated partial Euler product);
\item
  \(c_K(\chi,\rho) := \sum_{n \le K} \mu(n)\,\chi(n)\,n^{-\rho}\)
  (truncated Möbius sum);
\item
  \(D_K(\chi,\rho) := c_K(\chi,\rho)\,E_K(\chi,\rho)\)
  (Dirichlet \(D_K\) statistic);
\item
  \(T_K(\chi,\rho) := \sum_{p \le K}\sum_{k \ge 2} \chi(p)^k / (k\,p^{k\rho})\)
  and its limit \(T_\infty\) with \(B_\infty := \exp(T_\infty)\);
\item
  \(C_1(\chi,\rho) := -L''(\rho,\chi) / (2\,L'(\rho,\chi)^2)\)
  (subleading constant);
\item
  \(\psi\) for the primitive character of conductor \(f \mid q\) inducing
  \(\chi^2\).
\end{itemize}

The branch of \(\log L(2\rho,\psi)\) is fixed by analytic continuation
from \(\mathrm{Re}(s) > 1\), where the absolutely convergent log-Euler
expansion applies, to the boundary line \(\mathrm{Re}(s) = 1\) via
Hadamard--de la Vallée Poussin non-vanishing.

\subsection{Methodology of double verification}

Every numerical claim of \S\,X.5 is computed by \textbf{two independent
implementations in two languages with independent algorithms}.

\begin{itemize}
\item
  \textbf{L1 (primary).} \texttt{mpmath} 1.4 (Python 3.13), 50 decimal places.
  Direct partial-Möbius and partial-Euler evaluation; each zero \(\rho\)
  refined by Muller's method to \(|L(\rho,\chi)| < 10^{-50}\).
\item
  \textbf{L1b (in-language cross-check).} Same library, independent
  algorithm: Hurwitz-zeta expansion
  \(L(s,\chi) = q^{-s}\sum_{a=1}^{q}\chi(a)\,\zeta(s, a/q)\) in place
  of \texttt{mpmath.dirichlet}, central-difference numerical derivatives at
  three step sizes, and an independent linear sieve for \(\mu(n)\).
\item
  \textbf{L2 (cross-language).} PARI/GP 2.17.3 (C), 57 dps default; with
  python-flint 0.8.0 / Arb (FLINT 3.3) at 250 bits as a third stack
  on the worst-case pair.
\end{itemize}

Acceptance gates: agreement to \(\ge 12\) significant digits on \(\rho\)
and to \(\ge 8\) digits on \(L'\), \(L''\), \(C_1\); for residual quantities
where cancellation occurs (the \(B_\infty\) residual), agreement to
\(10^{-12}\) on each of the four component sums separately. Branch
choices and external citation provenance are recorded independently
in L1 and L2.

For the Phase-1 prime-residue replication of \S\,X.5.1, two analogous
stacks \textbf{P1a} and \textbf{P1b} are used: \texttt{primesieve} 12.13 (segmented
wheel sieve) and a hand-rolled plain-C segmented Eratosthenes sieve
with no external dependency.

\subsection{The local Perron double-pole residue}

\hypertarget{lemma:X.3.1}{\textbf{Lemma X.3.1.}} \emph{Let \(\chi\) be a primitive non-principal Dirichlet
character and let \(\rho\) be a simple zero of \(L(s,\chi)\). Then for
any \(K > 1\),}
\begin{equation}
\label{eq:res}
\mathop{\mathrm{Res}}_{w = 0}\!\left[\,\frac{K^{w}}{w\,L(w + \rho,\chi)}\,\right]
\;=\;
\frac{\log K}{L'(\rho,\chi)} \;+\; C_1(\chi,\rho).
\end{equation}

\emph{Proof.} Taylor-expand \(L(w + \rho,\chi)\) at \(w = 0\); invert to get
\(1/L(w+\rho,\chi) = (L'(\rho,\chi)\,w)^{-1} - L''(\rho,\chi)/(2\,L'(\rho,\chi)^2) + O(w)\).
Multiply by \(K^w/w = w^{-1} + \log K + O(w)\) and read off the
coefficient of \(w^{-1}\). \(\square\)

Lemma X.3.1 is unconditional given simplicity of \(\rho\). The full
algebraic derivation is in Appendix B \S\,B.2; the identity is also
machine-verified in Lean 4 / Mathlib v4.28.0 (\texttt{LocalPerronResidue.lean},
0 \texttt{sorry}; see \S\,X.6).

\subsection{Identities}

\subsubsection{Corrected \(B_\infty\) identity (unconditional)}

Our main algebraic result is the following unconditional identity
for the prime-power tail \(T_\infty\).

\hypertarget{theorem:X.4.1}{\textbf{Theorem X.4.1.}} \emph{Let \(\chi\) be a primitive non-principal Dirichlet
character of conductor \(q\), let \(\rho\) be a simple zero of
\(L(s,\chi)\) on the critical line, and let \(\psi\) be the primitive
character of conductor \(f \mid q\) inducing \(\chi^2\). Then}
\begin{equation}
\label{eq:Binfty}
T_\infty(\chi,\rho)
\;=\;
\tfrac12 \log L(2\rho,\psi)
\;+\; \mathrm{BPC}_1(\chi,\rho)
\;+\; \mathrm{BPC}_2(\chi,\rho)
\;+\; T_{\ge 3}(\chi,\rho),
\end{equation}
\emph{where}
\[
\mathrm{BPC}_1 = \tfrac12 \sum_{p \mid q,\, p \nmid f}
\log\bigl(1 - \psi(p)\,p^{-2\rho}\bigr),
\]
\[
\mathrm{BPC}_2 = -\tfrac12 \sum_{k \ge 2}\frac1k \sum_p \frac{\chi(p)^{2k}}{p^{2k\rho}},
\qquad
T_{\ge 3} = \sum_{k \ge 3}\frac1k \sum_p \frac{\chi(p)^k}{p^{k\rho}}.
\]
\emph{The four right-hand-side terms are individually finite. The identity
is unconditional given simplicity of \(\rho\).}

The \(k = 1\) prime sum \(\sum_p \chi^2(p) / p^{2\rho}\) is only
conditionally convergent on \(\mathrm{Re}(s) = 1\); its convergence is
supplied by Akatsuka~\cite[Lemma~2.1 \& eq.~(2.5)]{Akatsuka2017EulerProduct}, which is itself
unconditional (derived from PNT with explicit error term). \(\mathrm{BPC}_2\)
and \(T_{\ge 3}\) are absolutely convergent (minimum exponents
\(\mathrm{Re}(2k\rho) = 2\) and \(\mathrm{Re}(k\rho) = \tfrac32\)). The
full proof is given in Appendix A.

\subsubsection{Subleading constant \(C_1\) and the partial Möbius identity}

\hypertarget{theorem:X.4.2}{\textbf{Theorem X.4.2.}} \emph{Under the hypotheses of Theorem X.4.1, and
assuming every off-target nontrivial zero \(\rho'\) of \(L(s,\chi)\) with
\(|\mathrm{Im}(\rho')| \le T(K)\) is simple (for a zero-avoiding
truncation height \(T(K) \asymp K (\log K)^{-B}\) in the sense of
Inoue~\cite[Theorem~1]{Inoue2021}; we reserve \(T_K\) for the partial prime-power
sum of \S\,X.1),}
\begin{equation}
\label{eq:cK}
c_K(\chi,\rho) \;=\; \frac{\log K}{L'(\rho,\chi)} \;+\; C_1(\chi,\rho) \;+\; o(1)
\qquad (K \to \infty).
\end{equation}
\emph{The identity is unconditional in \(\rho\) given off-target simplicity.
The \(o(1)\) rate \(O(K^{-1/2+\epsilon})\) (every \(\epsilon > 0\)) and
the sharper \(O(K^{-1/2}\exp((\log K)^{1/2}(\log\log K)^{14}))\) via
the character analogue of Soundararajan~\cite[Theorem~1]{Soundararajan2009} are both
conditional on RH for \(L(s,\chi)\). For the four characters of \S\,X.5
the nontrivial zeros of \(L(s,\chi)\) in the relevant
explicit-formula range are numerically verified to lie on the
critical line (provenance in the Supplementary computation audit),
so these RH-conditional rates are the operative ones for the
finite \(K\) reported here; the unconditional fallback is the weaker
Akatsuka~\cite{Akatsuka2017EulerProduct} eq. (2.5) bound.}

The proof combines Inoue~\cite[Theorem~1]{Inoue2021}'s truncated explicit
formula for \(M^*(K,\chi)\) with Lemma X.3.1 to extract the double-pole
residue at \(\rho\); off-target zero contributions are bounded by
Soundararajan's argument. See Appendix B for the full proof.

\subsubsection{The Aoki--Koyama--Mertens constant (Hypothesis AK)}

Aoki--Koyama~\cite[eq.~(1.4), p.~235]{AokiKoyama2023}
states for a non-principal Dirichlet \(\chi\) that, under DRH,
\[
\lim_{x \to \infty}\Bigl((\log x)^{m}\!\!\prod_{p \le x}(1 - \chi(p)/p^s)^{-1}\Bigr)
\;=\;
\frac{L^{(m)}(s,\chi)}{e^{m\gamma}\,m!} \cdot
\begin{cases}\sqrt 2, & \chi^2 = 1,\ s = \tfrac12,\\ 1, & \text{otherwise,}\end{cases}
\]
with \(m = m(s,\chi) := \mathrm{ord}_{s' = s}\,L(s',\chi)\), the order
of zero of \(L(s,\chi)\) at the evaluation point \(s\) (so \(m\) is a
function of \(s\), not a fixed property of \(\chi\)). Specialised to the
regime of this paper (a simple noncentral zero, i.e.~\(m = 1\) at
\(s = \rho \ne \tfrac12\), branch multiplier \(1\)), this reads:
\begin{equation}
\label{eq:AK}
\lim_{K \to \infty} E_K(\chi,\rho)\,\log K \;=\; \frac{L'(\rho,\chi)}{e^{\gamma}}.
\tag*{(AK)}
\end{equation}
This \textbf{corrects} the earlier target \(L'(\rho,\chi) / \zeta(2)\): the
ratio is \(\zeta(2)/e^\gamma \approx 1.0828\). (AK) is DRH-conditional
in the form stated by Aoki--Koyama.

\subsubsection{The conditional NDC limit (open)}

Composing (\ref{eq:cK}) with (AK) formally gives
\begin{equation}
\label{eq:NDC}
D_K(\chi,\rho) = c_K(\chi,\rho)\,E_K(\chi,\rho)
\;\longrightarrow\; e^{-\gamma}
\qquad (K \to \infty),
\tag*{(NDC)}
\end{equation}
the corrected Numerical Duality Constant. The composition is
mechanical provided (\ref{eq:cK}) is strengthened to the \textbf{shifted
Perron leading theorem}:
\begin{equation}
\label{eq:Perron-leading}
c_K(\chi,\rho) \;=\; \frac{\log K}{L'(\rho,\chi)} \;+\; o(\log K)
\qquad (K \to \infty).
\tag*{(SP-L)}
\end{equation}
The obstruction is the off-target nontrivial-zero residue aggregate.
DRH constrains the location of off-target zeros to the critical line
but not their multiplicity; an off-target zero \(\lambda\) of
multiplicity \(m \ge 2\) contributes
\[
K^{\lambda - \rho}\,(\log K)^{m-1}\big/\bigl((m-1)!\,(\lambda - \rho)\,a_m(\lambda)\bigr),
\]
and the \((\log K)^{m-1}\) factor is not absorbed by simplicity of the
target \(\rho\). A literature search (\cite{Inoue2021}, \cite{Soundararajan2009},
Ng~\cite{Ng2004}, \cite{Akatsuka2017EulerProduct}) did not surface a theorem that closes
(SP-L); we state (NDC) as \textbf{conditional on (AK) and (SP-L)}
(Question Q:Perron, \S\,X.7).

\subsection{Numerical findings}

\subsubsection{Phase-1 Dominance-of-\(-1\) replication, \(x = 1.3 \cdot 10^{13}\)}

We independently replicate the prime-residue counts \(\pi(x; N, a)\)
of Koyama's \emph{nontriv.pdf} for \(N \in \{7, 8, 11, 19, 23\}\) and
\(x \in \{10^{12}, 1.3\cdot 10^{12}, 10^{13}, 1.3\cdot 10^{13}\}\) via
two independent prime-enumeration implementations (\texttt{primesieve} 12.13
and a hand-rolled segmented C sieve). Headline numbers:

\begin{itemize}
\item
  \(\pi(1.3 \cdot 10^{13}) = 445{,}831{,}610{,}611\), cross-checked
  against \texttt{primesieve\ -\/-count} standalone.
\item
  Library-independence: at every one of the four checkpoints, the
  \texttt{primesieve} counts and the hand-rolled C-sieve counts agree on
  every residue class for every \(N\).
\item
  Hardware-independence: a second M1-class machine agrees through
  \(x = 1.3 \cdot 10^{12}\) on every residue class for every \(N\).
\item
  Koyama's identity (3.1), a Dirichlet-orthogonality cross-check on
  the residue-count vector, is verified directly at all \(495\)
  \((N, x, a)\)-cells (worst absolute residual \(1.4 \cdot 10^{-4}\)).
  Here \(495\) counts \emph{every} residue class \(a\) across all five
  moduli and four checkpoints (the internal orthogonality identity
  holds for every class); the \(92\)-cell table below is the
  distinct subset of \((N, x, a)\) values that appear in Koyama's
  \emph{published} Tables 3--7 and can therefore be compared
  number-for-number against his manuscript.
\end{itemize}

Cell-by-cell comparison with Koyama's Tables 3--7 at all four
checkpoints, all moduli:

{\def\LTcaptype{none} % do not increment counter
\begin{longtable}[]{@{}
  >{\raggedright\arraybackslash}p{(\linewidth - 8\tabcolsep) * \real{0.1667}}
  >{\raggedleft\arraybackslash}p{(\linewidth - 8\tabcolsep) * \real{0.2222}}
  >{\raggedleft\arraybackslash}p{(\linewidth - 8\tabcolsep) * \real{0.2222}}
  >{\raggedleft\arraybackslash}p{(\linewidth - 8\tabcolsep) * \real{0.2222}}
  >{\raggedright\arraybackslash}p{(\linewidth - 8\tabcolsep) * \real{0.1667}}@{}}
\toprule\noalign{}
\begin{minipage}[b]{\linewidth}\raggedright
Table
\end{minipage} & \begin{minipage}[b]{\linewidth}\raggedleft
\(N\)
\end{minipage} & \begin{minipage}[b]{\linewidth}\raggedleft
cells
\end{minipage} & \begin{minipage}[b]{\linewidth}\raggedleft
exact
\end{minipage} & \begin{minipage}[b]{\linewidth}\raggedright
substantive disagreement (at \(x = 1.3 \cdot 10^{13}\))
\end{minipage} \\
\midrule\noalign{}
\endhead
\bottomrule\noalign{}
\endlastfoot
3 & 7 & 12 & 11 & 1 cell (\(\Delta = 50\), clean digit-shift profile) \\
4 & 8 & 12 & 1 & 11 (small-\(x\) rows; possible \(x\)-label error in Table 4 draft) \\
5 & 11 & 20 & 19 & 1 cell (\(a = 10\): our \(11{,}503\) vs Koyama \(71{,}711\)) \\
6 & 19 & 18 & 15 & 3 (2 substantive at \(a = 13, 18\); 1 sign flip at small \(x\)) \\
7 & 23 & 30 & 29 & 1 cell (\(\Delta = 100\), clean transposition profile) \\
\textbf{Total} & & \textbf{92} & \textbf{75} & \textbf{17} (\(74/81 \approx 91\%\) excluding the 11 Table-4 small-\(x\) rows) \\
\end{longtable}
}

Comparing to the qualitative dominance-of-\(-1\) statement of Koyama
(\emph{nontriv.pdf}, \S\,3): the signal is \textbf{cleanly reproduced for
\(N = 8\)} (\(-1 = 7\) is the strict largest at \(1.3 \cdot 10^{13}\),
\(\pi(\cdot ; 8, 7) - \pi(\cdot ; 8, 1) = 164{,}958\) vs \(126{,}732\)
and \(102{,}728\)). For \(N = 19\) the ranking of \(-1\) at this checkpoint
is 3rd of 9 non-residues in both our run and Koyama's table ---
consistent with \(-1\) being in the top group but not strictly
dominant. For \(N = 11\) the dominance turns on the disputed cell
\(a = 10\) (Table 5; see below); with Koyama's reported \(71{,}711\),
\(-1\) ranks 2nd of 5 (top group); with our \(11{,}503\), \(-1\) ranks
3rd of 5 (outside top group). Pending reconciliation of that cell.
For \(N = 7\) and \(N = 23\), Koyama himself notes (\emph{nontriv.pdf} p.~19)
that the predicted bias is not yet cleanly observed at this scale
and attributes it to the exceptionally low-lying first zero of the
relevant \(L\)-functions; our data agrees (\(-1\) is 2nd of 3 for
\(N = 7\), mid-rank for \(N = 23\)). Koyama's illustrative estimate
(\emph{nontriv.pdf} p.~19, for \(N = 19\)) places the next sine-wave
peak of the leading complex character at
\(x = e^{33.4} \approx 3.2 \cdot 10^{14}\), indicating that
strict dominance for the harder moduli is expected only at
substantially larger \(x\).

The full replication bundle (source code, build hashes, TSV
outputs, and reproducibility manifest) is deposited as
Supplementary Material S1.

\subsubsection{Numerical values of the four Dirichlet pairs}

The four pairs \((\chi, \rho)\) used throughout \S\,X.5.2--\S\,X.5.4 (all
values computed in mpmath at 50 decimal places):

{\def\LTcaptype{none} % do not increment counter
\begin{longtable}[]{@{}
  >{\raggedright\arraybackslash}p{(\linewidth - 8\tabcolsep) * \real{0.2000}}
  >{\raggedright\arraybackslash}p{(\linewidth - 8\tabcolsep) * \real{0.2000}}
  >{\raggedright\arraybackslash}p{(\linewidth - 8\tabcolsep) * \real{0.2000}}
  >{\raggedright\arraybackslash}p{(\linewidth - 8\tabcolsep) * \real{0.2000}}
  >{\raggedright\arraybackslash}p{(\linewidth - 8\tabcolsep) * \real{0.2000}}@{}}
\toprule\noalign{}
\begin{minipage}[b]{\linewidth}\raggedright
Pair
\end{minipage} & \begin{minipage}[b]{\linewidth}\raggedright
\(\chi\) (conductor)
\end{minipage} & \begin{minipage}[b]{\linewidth}\raggedright
\(\rho = \tfrac12 + i\tau\)
\end{minipage} & \begin{minipage}[b]{\linewidth}\raggedright
\(L'(\rho,\chi)\)
\end{minipage} & \begin{minipage}[b]{\linewidth}\raggedright
\(L''(\rho,\chi)\)
\end{minipage} \\
\midrule\noalign{}
\endhead
\bottomrule\noalign{}
\endlastfoot
\(\chi_{-4}/z_1\) & \(\chi_{-4}\) (\(q = 4\)) & \(\tfrac12 + 6.020949 i\) & \(\phantom{-}1.296500 + 0.182765 i\) & \(-1.697050 - 0.554017 i\) \\
\(\chi_{-4}/z_2\) & \(\chi_{-4}\) & \(\tfrac12 + 10.243770 i\) & \(\phantom{-}1.788467 - 0.296776 i\) & \(-3.319767 + 0.755548 i\) \\
\(\chi_5\) & \(\chi_5\) (\(q = 5\)) & \(\tfrac12 + 6.183578 i\) & \(\phantom{-}1.112930 - 0.448830 i\) & \(-1.642973 + 1.035107 i\) \\
\(\chi_{11}\) & \(\chi_{11}\) (\(q = 11\)) & \(\tfrac12 + 3.547041 i\) & \(\phantom{-}1.696582 - 0.250988 i\) & \(-3.121598 + 0.261219 i\) \\
\end{longtable}
}

L1b in-language cross-check (Hurwitz expansion, independent sieve)
agrees with L1 to \(|\Delta L'|, |\Delta L''| \lesssim 6 \cdot 10^{-12}\)
and \(|\Delta C_1| \lesssim 5 \cdot 10^{-13}\). L2 cross-language (PARI/GP
2.17.3) agrees with L1 to \(\ge 11\) decimal digits on every real and
imaginary component. An Arb spot-check at 250 bits on the worst pair
gives interval agreement on \(|L'|\) within \(3 \cdot 10^{-43}\).

\begin{quote}
\textbf{Reproducibility note (2026-05-16, D3 hardening).} The PARI/GP
2.17.3 (L2) and native 250-bit Arb cross-checks above were produced
in a prior environment and are \emph{not} re-runnable in the current one
(no \texttt{gp} binary; Arb here is python-flint 0.6.0). The hardened
verifier \texttt{handoff-2026-05-16-D3-binfty-hardening/binfty\_hardened.py}
supplies an \emph{independent, fully reproducible} substitute: two
engines --- \textbf{mpmath at dps 50 and 80} (precision-doubling) and
\textbf{python-flint / Arb 0.6.0} (rigorous ball arithmetic with proven
radii) --- agreeing to \(0\) at displayed precision on the
\(K\)-independent base \(\tfrac12\log L(2\rho,\psi)+\mathrm{BPC}_1+
\mathrm{BPC}_2\), with \(|L(\rho,\chi)|<10^{-67}\) at every refined
zero in both engines. It also reproduces the \(L',L'',C_1\) values
of the table above to all displayed digits. Referees should treat
the multi-engine evidence at exactly this reproducible strength;
the PARI/GP and native-Arb lines should be re-verified by the
author before submission or relabelled accordingly.
\end{quote}

\subsubsection{The Aoki--Koyama drift: \(e^{-\gamma}\) vs \(\zeta(2)^{-1}\)}

\emph{(Analytic-identity scale \(K \le 10^{7}\); not transferred from \S\,X.5.1.)}

The modulus \(|D_K|\) statistic at \(K = 2 \cdot 10^{6}\) and \(K = 10^{7}\)
(40 dps, mean over the four pairs):

{\def\LTcaptype{none} % do not increment counter
\begin{longtable}[]{@{}
  >{\raggedright\arraybackslash}p{(\linewidth - 8\tabcolsep) * \real{0.1579}}
  >{\raggedleft\arraybackslash}p{(\linewidth - 8\tabcolsep) * \real{0.2105}}
  >{\raggedleft\arraybackslash}p{(\linewidth - 8\tabcolsep) * \real{0.2105}}
  >{\raggedleft\arraybackslash}p{(\linewidth - 8\tabcolsep) * \real{0.2105}}
  >{\raggedleft\arraybackslash}p{(\linewidth - 8\tabcolsep) * \real{0.2105}}@{}}
\toprule\noalign{}
\begin{minipage}[b]{\linewidth}\raggedright
Quantity
\end{minipage} & \begin{minipage}[b]{\linewidth}\raggedleft
\(K = 2 \cdot 10^{6}\)
\end{minipage} & \begin{minipage}[b]{\linewidth}\raggedleft
\(K = 10^{7}\)
\end{minipage} & \begin{minipage}[b]{\linewidth}\raggedleft
\(\zeta(2)^{-1}\) target
\end{minipage} & \begin{minipage}[b]{\linewidth}\raggedleft
\(e^{-\gamma}\) target
\end{minipage} \\
\midrule\noalign{}
\endhead
\bottomrule\noalign{}
\endlastfoot
Mean \(|D_K| \cdot \zeta(2)\) & \(0.992\) & \(0.974\) & \(1.000\) & \(\zeta(2)\,e^{-\gamma} \approx 0.9237\) \\
Mean \(|E_K \log K|\,e^{\gamma}/|L'|\) & n/a & \(0.942\) & n/a & \(1.000\) \\
\end{longtable}
}

The drift from \(0.992\) to \(0.974\) between \(K = 2 \cdot 10^{6}\) and
\(K = 10^{7}\) is consistent with the AK normalisation \(e^{-\gamma}\)
at the natural \(1/\log K\) finite-size scale and \textbf{incompatible with
the \(\zeta(2)^{-1}\) target} at the same scale. We do not claim
convergence of the complex \(D_K(\chi,\rho)\) from a modulus statistic
alone --- that depends on (SP-L), which is open.

\subsubsection{The \(B_\infty\) identity at the four pairs}

\emph{(Analytic-identity scale \(K \le 10^{8}\); not transferred from \S\,X.5.1.)}

Identity residual \(|T_K - \mathrm{RHS}|\) for (\ref{eq:Binfty}) at
three scales --- \(K = 2 \cdot 10^{6}\) (mpmath, 50 dps; cross-checked
in PARI/GP 2.17.3, closed-form component evaluation), \(K = 10^{7}\)
and \(K = 10^{8}\) (PARI/GP):

{\def\LTcaptype{none} % do not increment counter
\begin{longtable}[]{@{}
  >{\raggedright\arraybackslash}p{(\linewidth - 8\tabcolsep) * \real{0.1579}}
  >{\raggedleft\arraybackslash}p{(\linewidth - 8\tabcolsep) * \real{0.2105}}
  >{\raggedleft\arraybackslash}p{(\linewidth - 8\tabcolsep) * \real{0.2105}}
  >{\raggedleft\arraybackslash}p{(\linewidth - 8\tabcolsep) * \real{0.2105}}
  >{\raggedleft\arraybackslash}p{(\linewidth - 8\tabcolsep) * \real{0.2105}}@{}}
\toprule\noalign{}
\begin{minipage}[b]{\linewidth}\raggedright
Pair
\end{minipage} & \begin{minipage}[b]{\linewidth}\raggedleft
\(K = 2 \cdot 10^{6}\) (L1)
\end{minipage} & \begin{minipage}[b]{\linewidth}\raggedleft
\(K = 2 \cdot 10^{6}\) (L2)
\end{minipage} & \begin{minipage}[b]{\linewidth}\raggedleft
\(K = 10^{7}\) (L2)
\end{minipage} & \begin{minipage}[b]{\linewidth}\raggedleft
\(K = 10^{8}\) (L2)
\end{minipage} \\
\midrule\noalign{}
\endhead
\bottomrule\noalign{}
\endlastfoot
\(\chi_{-4}/z_1\) & \(2.85 \cdot 10^{-3}\) & \(2.85 \cdot 10^{-3}\) & \(2.58 \cdot 10^{-3}\) & \(2.25 \cdot 10^{-3}\) \\
\(\chi_{-4}/z_2\) & \(1.66 \cdot 10^{-3}\) & \(1.66 \cdot 10^{-3}\) & \(1.52 \cdot 10^{-3}\) & \(1.32 \cdot 10^{-3}\) \\
\(\chi_5\) & \(4.24 \cdot 10^{-5}\) & \(4.24 \cdot 10^{-5}\) & \(1.22 \cdot 10^{-5}\) & \(3.30 \cdot 10^{-6}\) \\
\(\chi_{11}\) & \(3.34 \cdot 10^{-5}\) & \(3.34 \cdot 10^{-5}\) & \(1.75 \cdot 10^{-5}\) & \(4.10 \cdot 10^{-6}\) \\
\end{longtable}
}

L1 and L2 agree to all displayed digits at \(K = 2 \cdot 10^{6}\)
(stack difference \(\le 10^{-8}\)). \emph{(2026-05-16: the L2/PARI and
\(K=10^{7},10^{8}\) columns are from a prior environment and were not
re-run here; see the Reproducibility note in \S\,X.5.2. The hardened
verifier re-isolates the genuine analytic object --- the \(k=2\)
boundary identity \(R_2(K)=\tfrac12\sum_{p\le K}\chi^2(p)p^{-2\rho}
-[\tfrac12\log L(2\rho,\psi)+\mathrm{BPC}_1+\mathrm{BPC}_2]\),
removing the absolutely-convergent \(k\ge3\) tail that the L1/L2
\(|T_K-\mathrm{RHS}|\) figures conflate with it --- and confirms
\(R_2(K)\to0\) at the labelled rates across two engines.)} The decay across three decades
on the clean-character pairs \(\chi_5\), \(\chi_{11}\) (where
\(\chi(2) \ne 0\) and there is no bad-prime contribution to
\(\mathrm{BPC}_1\)):

{\def\LTcaptype{none} % do not increment counter
\begin{longtable}[]{@{}
  >{\raggedright\arraybackslash}p{(\linewidth - 8\tabcolsep) * \real{0.1579}}
  >{\raggedleft\arraybackslash}p{(\linewidth - 8\tabcolsep) * \real{0.2105}}
  >{\raggedleft\arraybackslash}p{(\linewidth - 8\tabcolsep) * \real{0.2105}}
  >{\raggedleft\arraybackslash}p{(\linewidth - 8\tabcolsep) * \real{0.2105}}
  >{\raggedleft\arraybackslash}p{(\linewidth - 8\tabcolsep) * \real{0.2105}}@{}}
\toprule\noalign{}
\begin{minipage}[b]{\linewidth}\raggedright
Pair
\end{minipage} & \begin{minipage}[b]{\linewidth}\raggedleft
residual ratio \(K = 2 \cdot 10^{6} \to 10^{7}\)
\end{minipage} & \begin{minipage}[b]{\linewidth}\raggedleft
\(K = 10^{7} \to 10^{8}\)
\end{minipage} & \begin{minipage}[b]{\linewidth}\raggedleft
\(\sqrt{5} \approx 2.24\)
\end{minipage} & \begin{minipage}[b]{\linewidth}\raggedleft
\(\sqrt{10} \approx 3.16\)
\end{minipage} \\
\midrule\noalign{}
\endhead
\bottomrule\noalign{}
\endlastfoot
\(\chi_5\) & 3.5 & 3.7 & 2.24 & 3.16 \\
\(\chi_{11}\) & 1.9 & 4.3 & 2.24 & 3.16 \\
\end{longtable}
}

Both pairs' ratios sit within a factor of \(\le 1.7\) of the
predicted \(K^{-1/2}\) rate. The relevant conditional input here is
RH for \(L(s, \chi^2)\) (equivalently \(L(s,\psi)\)) --- the \(B_\infty\)
residual's slow component is the boundary-line \(k = 1\) sum
\(\sum_p \chi^2(p)\,p^{-2\rho}\), governed by the \(\chi^2\)/\(\psi\)
\(L\)-function; this is a \emph{different} \(L\)-function from the
\(L(s,\chi)\) that governs the \(c_K\) rate of Theorem X.4.2, because
it is a different partial sum. The square-root-type decay is the
character analogue of Soundararajan~\cite{Soundararajan2009}'s RH-conditional bound.
We do not claim this rate unconditionally: what is unconditional
is the (much weaker) \cite{Akatsuka2017EulerProduct} eq. (2.5) bound \(O(1/\log K)\)
for that \(k = 1\) partial sum. The observed decay is faster than
this unconditional floor and consistent with the RH-conditional
rate; for \(\chi_{-4}, \chi_5, \chi_{11}\) the relevant zeros of
\(L(s,\chi^2)\) lie on the critical line throughout the numerically
verified range (provenance in the Supplementary computation
audit), so the RH-conditional rate is the operative one across
the \(K\)-scales reported here.
\(\chi_5\) sits consistently above the prediction, \(\chi_{11}\)
straddles it across the two \(K\)-steps --- well within the
oscillatory \(O(1)\) implicit-constant envelope of the
Soundararajan--conditional rate. The \(\chi_{-4}\) pairs show systematically slower decay
(ratio \(1.09\)--\(1.15\) across the two K-steps), consistent with the
additional bad-prime \(p = 2\) contribution to \(\mathrm{BPC}_1\)
suppressing the leading-order convergence rate.

\subsubsection{Two negative elliptic-curve findings}

We record two negative findings on the elliptic-curve side, both
deferred to the supplementary record for the design data, source,
and resampling diagnostics:

\begin{itemize}
\item
  \emph{Conductor-confounded rank trend.} A multivariate OLS regression
  of \(E[C_1^2]\) on \((\mathrm{rank}, \log N)\) across 19 weight-\(2\)
  elliptic curves shows a stable \(\log N\) coefficient and an
  \emph{unstable} rank coefficient (bootstrap \(95\%\) CI for the rank
  coefficient includes zero; one rank-\(3\) anchor swings the
  coefficient by \(63\%\)). We describe this as a
  conductor-confounded trend, not a rank law (Question Q:conductor,
  \S\,X.7).
\item
  \emph{Raw \(\mathrm{Sym}^2 / \langle f, f \rangle\) proportionality.}
  The raw ratio ranges over seven orders of magnitude across
  \(\{37a_1, 389a_1, \Delta\}\) and is empirically falsified in its
  raw form (Question Q:Sym2, \S\,X.7).
\end{itemize}

\subsection{Lean 4 / Mathlib4 formalisation}

We accompany the paper with a Lean 4 / Mathlib4 lake project
(\texttt{formal-conjectures/} directory of the supplementary archive).
The toolchain is \texttt{leanprover/lean4:v4.28.0}; Mathlib is pinned at
commit \texttt{8f9d9cff6bd728b17a24e163c9402775d9e6a365}. The Lean
inventory fixes
the \textbf{statements} of every identity of \S\,X.4, ensures normalisations
and branch conventions are syntactically explicit, and records each
statement's proof status against a public audit trail.

\textbf{Build status.} The \S\,X formalisation comprises the \textbf{10 Lean
modules} enumerated in the inventory table below. The lake
roll-up target \texttt{FormalConjectures} builds these together with one
further module, \texttt{SignedVsAbsoluteResidueGadget.lean} (a halo-route
structural lemma belonging to the companion GL(2) strand of
Question Q:EC-recip, \S\,X.7, and outside the scope of \S\,X), for a
total of \textbf{11 modules in the build target}, all of which compile
under \texttt{leanprover/lean4:v4.28.0}. (The \texttt{formal-conjectures/}
directory additionally holds the transient \texttt{\_AxiomCheck.lean}
audit harness and one round-9 scratch extract, neither part of the
build target; file counts in the directory listing should not be
read as module counts.) Across the build there are exactly
\textbf{2 \texttt{sorry}s}, both the DPAC headline conjecture at general \(K\):
one in \texttt{DPAC\_full.lean:338} (the obstruction annotated in-source
as \texttt{-\/-\ RESEARCH-OPEN:} at line 321) and one in
\texttt{DirichletPolynomialAvoidance.lean:54}, a statement-only mirror of
the conjecture (Saar Shai, \emph{Prime Spectroscopy of Riemann Zeros},
\S\,3) carrying a bare \texttt{sorry} with no in-source annotation or
category attribute. Of the 10 \S\,X modules, \textbf{8 are fully proved
(0 \texttt{sorry})} --- and the out-of-scope 11th module is likewise
\texttt{sorry}-free --- covering the algebraic content of \S\,X.3, \S\,X.4, the
smoothed explicit-formula chain, the Mertens spectroscope
universality statement, the Farey bridge identity (the
underlying static Farey--Mertens identity is classical, in the
Mikol\'as~\cite{Mikolas1949} tradition; what is contributed here is its
unconditional Lean formalisation, now that \texttt{RamanujanSum.lean}
has discharged the Ramanujan-sum hypothesis), the Farey
sign-pattern statement
(conditional on \texttt{h\_chebyshev\_bias} and the two pointwise
falsification witnesses at \(p = 237{,}733\) and \(p = 243{,}799\)),
and DPAC for \(K \le 4\).

No \texttt{axiom} declarations are introduced anywhere in the project. A
companion \texttt{\_AxiomCheck.lean} file runs \texttt{\#print\ axioms} on each
audited headline theorem. Six audited headlines (the \texttt{RamanujanSum}
chain, \texttt{FareyBridgeIdentity}, \texttt{LocalPerronResidue},
\texttt{CorrectedBInfty}, \texttt{MertensSpectroscopeUniversality},
\texttt{FareySignPattern}) depend only on the standard Lean trust base
\texttt{propext}, \texttt{Classical.choice}, \texttt{Quot.sound}. The remaining audited
headline \texttt{dpac\_le\_4} (unconditional DPAC for \(K \in \{2, 3, 4\}\))
additionally uses \texttt{Lean.ofReduceBool} and \texttt{Lean.trustCompiler} ---
Mathlib's standard kernel-reduction primitives, used because the
proof evaluates Möbius values at small primes in the kernel. The
eighth fully-proved file, \texttt{SmoothedDwfFormula\_full}, is a 17-lemma
algebraic-glue chain rather than a single headline theorem; its
component lemmas use only the standard trust base, and its two
analytic prerequisites are stated as explicit hypotheses on the
consuming theorems. No axiom is unstable or project-specific. The
full per-\texttt{sorry} inventory and the cumulative axiom audit are in
the companion \texttt{LEAN\_SORRY\_STATUS.md} of the reproducibility bundle.

\textbf{Status of headline theorems} in \S\,X.3--\S\,X.4:

{\def\LTcaptype{none} % do not increment counter
\begin{longtable}[]{@{}
  >{\raggedright\arraybackslash}p{(\linewidth - 2\tabcolsep) * \real{0.5000}}
  >{\raggedright\arraybackslash}p{(\linewidth - 2\tabcolsep) * \real{0.5000}}@{}}
\toprule\noalign{}
\begin{minipage}[b]{\linewidth}\raggedright
Paper object
\end{minipage} & \begin{minipage}[b]{\linewidth}\raggedright
Lean status
\end{minipage} \\
\midrule\noalign{}
\endhead
\bottomrule\noalign{}
\endlastfoot
Lemma X.3.1 (local Perron residue) & \textbf{Lean-verified unconditional} (\texttt{LocalPerronResidue.lean}, 0 \texttt{sorry}) \\
Theorem X.4.1 (corrected \(B_\infty\) identity) & \textbf{Lean-verified conditional} on the convergence-of-partial-sum hypothesis derived in Appendix A (\texttt{CorrectedBInfty.lean}, 0 \texttt{sorry}) \\
Theorem X.4.2 (\(c_K\) leading + subleading) & Pen-and-paper proof in Appendix B; off-target-zero-simplicity hypothesis stated explicitly \\
(AK), (SP-L), (NDC) & (AK) cited (Aoki--Koyama~\cite{AokiKoyama2023}); (SP-L) and (NDC) open challenges (Q:Perron, \S\,X.7) \\
DPAC & \textbf{Lean-verified for \(K \le 4\) unconditional} (\texttt{DPAC\_closure\_attempt.lean}, 0 \texttt{sorry}); general \(K\) open (LI-class), with obstruction certificate via \cite{Polya1913} \\
\end{longtable}
}

{\def\LTcaptype{none} % do not increment counter
\begin{longtable}[]{@{}
  >{\raggedright\arraybackslash}p{(\linewidth - 4\tabcolsep) * \real{0.3333}}
  >{\raggedright\arraybackslash}p{(\linewidth - 4\tabcolsep) * \real{0.3333}}
  >{\raggedright\arraybackslash}p{(\linewidth - 4\tabcolsep) * \real{0.3333}}@{}}
\toprule\noalign{}
\begin{minipage}[b]{\linewidth}\raggedright
Paper object
\end{minipage} & \begin{minipage}[b]{\linewidth}\raggedright
Lean file
\end{minipage} & \begin{minipage}[b]{\linewidth}\raggedright
Status
\end{minipage} \\
\midrule\noalign{}
\endhead
\bottomrule\noalign{}
\endlastfoot
Boundary residue \(R_0 = -2\) for a Gaussian-cutoff Mellin-shift explicit formula (companion strand) and its algebraic-glue chain & \texttt{SmoothedDwfFormula\_full.lean} & \textbf{THEOREM (chain), 0 \texttt{sorry}.} All 17 algebraic-glue lemmas closed unconditionally; the two analytic prerequisites \texttt{mellin\_decay} (Stirling on \(\Gamma\) vertical strips) and \texttt{inv\_zeta\_polynomial\_growth} (Titchmarsh~\cite[\S\,3.11]{Titchmarsh1986RZF}) are now stated as explicit hypotheses on the theorems that consume them, both Mathlib v4.28.0 gaps. \\
Lemma X.3.1 (local Perron residue) & \texttt{LocalPerronResidue.lean} & \textbf{THEOREM (0 \texttt{sorry}).} The residue identity is stated as a \texttt{Tendsto} limit at \(L\) analytic with simple zero at \(0\) (the general-\(\rho\) case reduces by \(L \mapsto L(\cdot + \rho)\)). \\
Theorem X.4.1 (\(B_\infty\) identity) & \texttt{CorrectedBInfty.lean} & \textbf{THEOREM (0 \texttt{sorry}), conditional on \texttt{h\_convergence}.} The four-component identity is proved against \texttt{noncomputable\ def}s of \(T_\infty\), \(T_{\ge 3}\), \(\mathrm{BPC}_1\), \(\mathrm{BPC}_2\), \(L\) given an added hypothesis \texttt{h\_convergence\ :\ Tendsto\ T\_K\ atTop\ (nhds\ (RHS))}. This hypothesis packages exactly the four analytic inputs of the pen-and-paper proof in Appendix A (\cite{Akatsuka2017EulerProduct}, log-Euler-product, imprimitive induction, geometric tails); the Lean proof uses \texttt{Classical.epsilon\_spec} + \texttt{tendsto\_nhds\_unique} and is three lines. \\
Farey bridge identity & \texttt{FareyBridgeIdentity.lean} & \textbf{THEOREM (0 \texttt{sorry}), unconditional} (\texttt{farey\_bridge\_identity\_unconditional}). The \texttt{h\_ramanujan\_decomp} hypothesis is now discharged by \texttt{RamanujanSum.farey\_ramanujan\_decomp}; the only inputs are \texttt{Nat.Prime\ p} and Mathlib v4.28.0. \emph{Provenance:} the underlying static Farey--Mertens identity (the \(m = p\) slice \(\sum_{f \in \mathcal{F}_{p-1}} e^{2\pi i p f} = M(p) + 2\)) is classical, in the Mikol\'as~\cite{Mikolas1949} tradition and a special case of the Farey-discrepancy Fourier spectrum; the contribution recorded here is the unconditional machine formalisation, not the identity itself. The genuinely novel mathematical content of this strand is the \emph{differential, per-step} refinement (the sign behaviour as a single prime denominator enters), developed in the companion Dominance/Chebyshev-bias chapter, not \S\,X. \\
Mertens spectroscope universality & \texttt{MertensSpectroscopeUniversality.lean} & \textbf{THEOREM (0 \texttt{sorry}), conditional on an explicit-formula asymptotic hypothesis} (Soundararajan~\cite[Theorem~1]{Soundararajan2009} input). The file additionally contains a 5-step blueprint documenting the precise Mathlib gap (Perron inversion, explicit formula for \(M(x)\), oscillatory-integral partial summation, zero simplicity) and two new unconditionally-proved infrastructure lemmas: \texttt{spectroscope\_nonneg} (the spectroscope statistic is non-negative) and \texttt{reciprocal\_sqrt\_not\_summable} (if \(\sum_{p \in P} 1/p\) diverges, so does \(\sum_{p \in P} 1/\sqrt p\)). \\
Farey sign pattern & \texttt{FareySignPattern.lean} & \textbf{THEOREM (0 \texttt{sorry}), conditional.} Three theorems closed under explicit named hypotheses: \texttt{farey\_sign\_pattern\_density\_one} (density-one \texttt{Tendsto} version, under \texttt{h\_chebyshev\_bias}); \texttt{pointwise\_falsification\_237733} and \texttt{pointwise\_falsification\_243799} (the two pointwise falsifications, each under a \texttt{h\_witness} hypothesis stating that the relevant signs disagree), packaged into \texttt{pointwise\_version\_falsified}. Negative result for the pointwise conjecture; the falsifications are recorded as theorems, not axioms --- the project's ``no \texttt{axiom}'' convention is preserved. \\
Ramanujan sum + Farey decomposition (Hardy \& Wright Thms 271, 304) & \texttt{RamanujanSum.lean} & \textbf{THEOREMS (0 \texttt{sorry}), unconditional.} Geometric sum identity for roots of unity (\texttt{geom\_sum\_roots\_of\_unity}); primitive-roots-sum equals Möbius (\texttt{primRootsSum\_eq\_moebius}, via Dirichlet convolution + strong induction); the coprime case \(c_q(n) = \mu(q)\) (\texttt{ramanujanSum\_eq\_moebius\_of\_coprime}); FareySet sum decomposition (\texttt{farey\_ramanujan\_decomp}) discharging the \texttt{h\_ramanujan\_decomp} hypothesis above. \\
Dirichlet Polynomial Avoidance (DPAC) --- partial closure + bridges & \texttt{DPAC\_closure\_attempt.lean}, \texttt{DirichletPolynomialAvoidance.lean}, \texttt{DPAC\_full.lean} & \textbf{PARTIAL.} \texttt{DPAC\_closure\_attempt.lean} (0 \texttt{sorry}) proves DPAC unconditionally for \(K \in \{2, 3, 4\}\) using only \(0 < \mathrm{Re}(\rho) < 1\) (\texttt{dpac\_K\_eq\_2}, \texttt{dpac\_K\_eq\_3}, \texttt{dpac\_K\_eq\_4}, \texttt{dpac\_le\_4}). It also reformulates the open case as \texttt{FiniteLogRatioLI} and records the obstruction certificate (\cite{Polya1913} discreteness of the exponential-polynomial zero set + a single open avoidance statement). The headline conjecture for general \(K\) remains \texttt{sorry} in \texttt{DPAC\_full.lean:338} and \texttt{DirichletPolynomialAvoidance.lean:54}, diagnostically comparable to the Linear Independence Hypothesis for \(\zeta\)-zero ordinates (made precise in the \S\,X.7 Structural remark); the four explicit phase-avoidance bridges (\texttt{dpac\_of\_logPrimePhaseAvoidance} through \texttt{dpac\_of\_certifiedZetaZeroSample}) are closed without \texttt{sorry}. \\
\end{longtable}
}

The role of the Lean artifact is to fix the statements and provide
a publicly inspectable audit trail of the proof obligations
remaining.

\subsection{Open challenges}

The following structure the next phase of the program.

\begin{quote}
\textbf{Q:Perron (Shifted Perron leading theorem).} Prove the shifted
Perron leading statement (SP-L) for primitive non-principal
\(\chi\) and simple non-central \(\rho\).
\end{quote}

Three sufficient packages, in decreasing strength of input
required (see \texttt{SP\_L\_SUFFICIENT\_PACKAGES\_2026-05-13.md} of the
supplementary archive for the full discussion):

\begin{itemize}
\item
  \textbf{Route I (off-target simplicity + shifted negative second
  moment).} All off-target zeros at the truncation height are
  simple, and
  \(\sum_{\rho'}^\mathrm{mult} |L(\rho' + \alpha,\chi)|^{-2} \ll_\chi (\log K)^{O(1)}\).
  Near-Lindelöf strength; not unconditionally available.
\item
  \textbf{Route II (halo-route reduction).} The cluster-summed-residue
  contour pivot (transferred from the GL(2) version) replaces the
  termwise budget by a signed cancellation but yields only
  \(|R_K| \ll K^{1/2 + \varepsilon}\) for the off-target aggregate ---
  \textbf{far above the \(o(\log K)\) target}. The halo route in its
  present form does not close (SP-L), though it does replace the
  rooted Palm wall as the structural obstruction.
\item
  \textbf{Route III (direct partial summation).} Substantially weaker:
  a Gonek--Hejhal-type bound
  \(A(T) := \sum_{\gamma' \le T,\, \gamma' \ne \tau}
  1/[(\rho' - \rho)\,|L'(\rho', \chi)|] \ll_\chi \log T\), combined
  with a Mertens-style oscillation bound
  \(\int_0^T K^{i(\gamma - \tau)}\, dA(\gamma) = o(\log T)\) uniformly
  in \(K\) along the zero-avoiding heights. This is the cleanest
  pen-and-paper analogue of \cite{Akatsuka2017EulerProduct} eq.\textasciitilde(2.5) applied to the
  off-target weight sequence.
\end{itemize}

The total-Möbius bounds of Soundararajan type are too coarse to
isolate the pointwise cancellation at the \(\log K\) scale by
themselves.

\begin{quote}
\textbf{Q:EC-recip (GL(2) reciprocal-derivative control).} Prove a
fixed-curve theorem for
\(\sum_\gamma \widehat W(i\gamma)\,e^{i\gamma u} / L'(E, 1 + i\gamma)\)
giving cancellation \(o(u^r)\), or a minimum-modulus estimate
on a vertical line with explicit exponent \(< 2\).
\end{quote}

Without a GL(2) analogue of Aoki--Koyama, the EC side remains at the
level of quantitative ensemble evidence.

\begin{quote}
\textbf{Q:DPAC (Dirichlet Polynomial Avoidance).} Prove DPAC for
general \(K\). We give an unconditional Lean proof for
\(K \in \{2, 3, 4\}\); the general case reduces, via \cite{Polya1913}
discreteness of the finite-exponential-polynomial zero set, to
a single open avoidance statement at \(\zeta\)-zero ordinates
(made precise in the Structural remark below).
\end{quote}

\textbf{Structural remark (shared obstruction).} We record an
identification of the \emph{form} of the obstruction common to the open
items above --- an observation about structure, not a resolution,
and not asserted as a theorem. The shifted-Perron leading
statement (SP-L) of Q:Perron and the general-\(K\) reduction
\texttt{FiniteLogRatioLI} of Q:DPAC are, respectively, the sharp
(\(c \to 1\)) and discrete instantiations of a single
negative-second-moment / linear-independence phenomenon for the
zero ordinates: both are controlled by a Gonek--Hejhal-type
reciprocal-derivative second moment of the
\(\sum_{\rho} |\zeta'(\rho)|^{-2}\) family (Ng~\cite{Ng2004}) together with a
quantitative linear-independence input on the relevant ordinates.
The companion GL(2) strand of Q:EC-recip reduces (softly, \(c < 3\))
to the same family at GL(2). Thus ``diagnostically comparable to
the Linear Independence Hypothesis'', used loosely above, is here
made precise: under a quantitative LI hypothesis for the relevant
\(L\)-function's zero ordinates the shared second moment is
controlled, simultaneously closing the GL(1) \texttt{FiniteLogRatioLI}
obstruction and feeding (SP-L); neither instantiation is proved
unconditionally, and locating this common barrier precisely (not
moving it) is the present contribution.

\textbf{Further questions} (from the EC and ensemble negatives of \S\,X.5.5
and from the \S\,X.5-companion EC-NDC programme; deferred to the
supplementary record):

\begin{itemize}
\item
  \emph{Q:conductor} --- Replicate the \S\,X.5.5 regression of
  \(\mathbb{E}[C_1^2]\) on \((\mathrm{rank}, \log N)\) on a curve set
  where rank and \(\log N\) are not collinear, to separate any
  genuine rank dependence from the conductor contribution.
\item
  \emph{Q:Sym2} --- Identify a completed or archimedean-corrected
  \(\mathrm{Sym}^2\) normalisation replacing the empirically falsified
  raw \(\mathrm{Sym}^2 / \langle f, f \rangle\) proportionality (which
  ranges over seven orders of magnitude across
  \(\{37a_1, 389a_1, \Delta\}\)).
\item
  \emph{Q:EC-NDC} --- Identify a normalisation of \(D_K^E\) for which the
  universal limit exists and is distinguishable from null
  transforms. The sharp-cutoff form \(D_K^E \cdot \zeta(2) \to 1\) is
  falsified through \(K = 10^{6}\); smoothed variants show numerical
  agreement, but the predeclared G3 specificity gate fails to
  separate them from null controls, so the apparent agreement is
  not yet significant.
\end{itemize}

\subsection{Code, data, and certificate availability}

All scripts, refined zero data, numerical-table CSVs, convergence
logs, the Lean 4 lake project, and the reproducibility manifest will
be deposited as a single self-contained reproducibility bundle
(Supplementary material S1), mirrored at a Zenodo DOI at acceptance.
The bundle pins all software versions: Lean toolchain
\texttt{leanprover/lean4:v4.28.0}, Mathlib commit
\texttt{8f9d9cff6bd728b17a24e163c9402775d9e6a365}, \texttt{mpmath} 1.4,
PARI/GP 2.17.3, FLINT 3.3 / python-flint 0.8.0. The Phase-1
Dominance-of-\(-1\) replication bundle is included as the
Supplementary Material S1 archive. Each numerical
table in \S\,X.5 cites the L1 script and L2 reproducer; each external
theorem cited in \S\,X.4 has its PDF retrieval recipe, page/equation,
and verbatim quote recorded in the citation audit
(Supplementary S2, \emph{Citation audit}).



\appendix
\input{appendix_A}
\section{Full proof of Theorem X.4.2 ($c_K$ leading + subleading identity)}

This appendix gives the full proof of Theorem X.4.2 of the main text:
under the hypotheses of Theorem X.4.1 (i.e., \(\chi\) primitive
non-principal of conductor \(q\), \(\rho = \tfrac12 + i\tau\) a \emph{simple}
zero of \(L(s,\chi)\) with \(\tau \ne 0\)), and assuming every
off-target nontrivial zero \(\rho' \ne \rho\) of \(L(s,\chi)\) in the
truncation height \(|\mathrm{Im}(\rho')| \le T(K)\) of the Inoue~\cite{Inoue2021}
explicit-formula box is also \emph{simple} (we write \(T(K)\) for the
truncation height, distinct from the partial prime-power sum
\(T_K(\chi,\rho)\) of \S\,X.1 of the main text; the \(o(1)\) rate further
uses RH for \(L(s,\chi)\), see \S\,B.3 and \S\,B.4),
\begin{equation}
\label{eq:cK-appendix}
c_K(\chi,\rho) \;=\; \frac{\log K}{L'(\rho,\chi)} \;+\; C_1(\chi,\rho) \;+\; o(1)
\qquad (K \to \infty),
\tag{$\dagger$}
\end{equation}
where \(C_1(\chi,\rho) = -L''(\rho,\chi) / (2\,L'(\rho,\chi)^2)\) and
\(c_K(\chi,\rho) = \sum_{n \le K} \mu(n)\,\chi(n)\,n^{-\rho}\).

The identity (\(\dagger\)) is \textbf{unconditional in \(\rho\)} given
simplicity of \(\rho\) and of the crossed off-target zeros. The \(o(1)\) rate
\(O(K^{-1/2 + \varepsilon})\) for every \(\varepsilon > 0\), and the
sharper rate
\(O\bigl(K^{-1/2} \exp\bigl((\log K)^{1/2}(\log\log K)^{14}\bigr)\bigr)\)
obtained via the character analogue of Soundararajan~\cite[Theorem~1]{Soundararajan2009},
are both conditional on RH for \(L(s, \chi)\). For the four characters
\(\chi_{-4}, \chi_5, \chi_{11}\) of \S\,X.5, the nontrivial zeros of
\(L(s,\chi)\) in the relevant explicit-formula range are numerically
verified to lie on the critical line (provenance in the Supplementary
computation audit), so these RH-conditional rates are the operative
ones for the finite \(K\) reported here; we do not assert them as
unconditional theorems. The unconditional fallback is the weaker
Akatsuka~\cite{Akatsuka2017EulerProduct} eq. (2.5) \(O(1/\log K)\) bound.

\subsection{Setup --- Perron formula for \(c_K\) via \(1/L(s, \chi)\)}

\subsubsection{The Dirichlet series for \(1/L(s, \chi)\)}

For \(\chi\) primitive of conductor \(q \ge 2\) and \(\mathrm{Re}(s) > 1\),
\[
\frac{1}{L(s, \chi)}
\;=\;
\sum_{n=1}^{\infty} \frac{\mu(n)\,\chi(n)}{n^s},
\]

absolutely convergent on \(\mathrm{Re}(s) > 1\). The Dirichlet
coefficients are \(a_n = \mu(n)\,\chi(n)\), and the partial sum
\(c_K(\chi, \rho) = \sum_{n \le K} a_n n^{-\rho}\) is the spectroscope
object (argument order matching \S\,X.1 of the main text).

\subsubsection{Perron formula (truncated)}

The standard truncated Perron formula (Tenenbaum, \emph{Introduction to
Analytic and Probabilistic Number Theory}, Theorem II.2.3; Inoue
(2021), \emph{Some explicit formulas for partial sums of Möbius
functions}, Journal de Théorie des Nombres de Bordeaux \textbf{33} (2021),
273--315, eq. (4.1) and Theorem 1) gives the following representation.

For any \(\sigma_0' > \tfrac12\) and \(T \ge 2\), with the convention that
\(c_K\) uses the half-weight at integer \(K\),
\begin{equation}
\label{eq:Perron-truncated}
c_K(\chi, \rho)
\;=\;
\frac{1}{2\pi i} \int_{\sigma_0' - iT}^{\sigma_0' + iT}
\frac{K^w}{w} \cdot \frac{1}{L(w + \rho, \chi)} \,dw
\;+\; \mathcal{R}'(K, T, \rho),
\end{equation}
where the kernel is \(K^w / w\) and the Dirichlet series factor is
\(1/L(w + \rho, \chi)\), obtained by the substitution \(s = w + \rho\).
The truncation error \(\mathcal{R}'(K, T, \rho)\) is bounded by Inoue
(2021) Theorem 1 via the \(J_1, J_2, J_3\) estimates of that paper:
\[
\mathcal{R}'(K, T, \rho)
\;\ll\;
\frac{K^{\sigma_0' - 1/2}\log K}{T}
\;+\;
\min\!\Bigl(1,\, \frac{K^{1/2}}{T\langle K\rangle}\Bigr),
\]

with \(\langle K \rangle\) the distance from \(K\) to the nearest integer.

\subsubsection{Inoue~\cite[Theorem~1]{Inoue2021} --- truncated explicit formula}

The following statement is verbatim from Inoue~\cite{Inoue2021}, arXiv:1805.05015v1, page 3:

\begin{quote}
\textbf{Theorem 1} (\cite{Inoue2021}). Let \(x > 0\), \(q \ge 2\),
\(T \ge \max\{T_0, \exp(q^{1/3}), 2/x\}\). Then, uniformly for all
primitive Dirichlet characters \(\chi\) modulo \(d\) with \(d \le q\),
there exists \(T_\nu \in [T, 2T]\) satisfying

\[M^*(x, \chi)
  \;=\; \sum_{|\gamma| < T_\nu} \frac{1}{(m(\rho)-1)!}\,
       \lim_{s \to \rho} \frac{d^{m(\rho)-1}}{ds^{m(\rho)-1}}
       \Bigl((s - \rho)^{m(\rho)}\,\frac{x^s}{L(s, \chi)\,s}\Bigr)
       \;+\; \mathop{\mathrm{Res}}_{s=0}\Bigl(\frac{x^s}{L(s, \chi)\,s}\Bigr)
       \;+\; \cdots\]
\end{quote}

This is the unshifted form, applied to the partial-character-Möbius
sum \(M^*(x, \chi) = \sum_{n \le x}' \mu(n)\,\chi(n)\). Dividing by
\(K^\rho\) and specializing at \(x = K\), \(m(\rho) = 1\) (the simplicity
hypothesis on \(\rho\)), and after the change of variable \(w = s - \rho\),
one obtains (\ref{eq:Perron-truncated}).

\subsection{Pole-structure analysis at \(w = 0\)}

The pole-structure computation of this section is independently
formalised in Lean 4 / Mathlib v4.28.0 as \texttt{local\_perron\_residue} in
\texttt{formal-conjectures/LocalPerronResidue.lean}. The Lean theorem
states the residue identity as a \texttt{Tendsto} limit at the canonical
form \(L\) analytic with a simple zero at \(0\) (the general-\(\rho\) case
of this appendix reduces to the canonical form by the substitution
\(L \mapsto L(\,\cdot\, + \rho)\)); the proof is \textbf{unconditional}
(0 \texttt{sorry}, only \texttt{propext}, \texttt{Classical.choice}, \texttt{Quot.sound}
axioms).

We work in the shifted variable \(w = s - \rho\). Define
\[
F(w) \;:=\; \frac{K^w}{w \cdot L(w + \rho, \chi)}.
\]

Since \(L(w + \rho, \chi)\) has a \emph{simple} zero at \(w = 0\) (by hypothesis),
\(1/L(w+\rho, \chi)\) has a simple pole at \(w = 0\). Combined with the
simple pole of \(K^w / w\) at \(w = 0\), the integrand \(F\) has a
\textbf{double pole at \(w = 0\)}.

\subsubsection{Laurent expansion of \(1/L(w+\rho, \chi)\)}

By Taylor expansion at the simple zero,
\[
L(w + \rho, \chi)
\;=\;
a_1 w + a_2 w^2 + a_3 w^3 + \cdots,
\qquad a_k \;=\; \frac{L^{(k)}(\rho, \chi)}{k!},\quad a_1 = L'(\rho, \chi) \ne 0.
\]

Therefore
\[
\frac{1}{L(w + \rho, \chi)}
\;=\;
\frac{1}{a_1 w} \cdot \frac{1}{1 + (a_2/a_1) w + (a_3/a_1) w^2 + \cdots}.
\]

Expanding the geometric factor as \((1 + u)^{-1} = 1 - u + u^2 - \cdots\)
and substituting \(a_1 = L'(\rho, \chi)\), \(a_2 = L''(\rho, \chi)/2\):
\begin{equation}
\label{eq:invL-Laurent}
\frac{1}{L(w + \rho, \chi)}
\;=\;
\frac{1}{L'(\rho, \chi)\,w}
\;-\;
\frac{L''(\rho, \chi)}{2\,L'(\rho, \chi)^2}
\;+\;
O(w).
\end{equation}

\subsubsection{Laurent expansion of \(K^w / w\)}

Trivially,
\[
\frac{K^w}{w}
\;=\;
\frac{1}{w}\bigl(1 + (\log K)\,w + \tfrac{1}{2}(\log K)^2 w^2 + O(w^3)\bigr)
\;=\;
\frac{1}{w} \;+\; \log K \;+\; \tfrac{1}{2}(\log K)^2 \,w \;+\; O(w^2).
\]

\subsubsection{Product expansion and the double-pole residue}

Multiplying the two expansions:
\[
F(w)
\;=\;
\Bigl[\tfrac{1}{w} + \log K + \tfrac{1}{2}(\log K)^2 w + \cdots\Bigr]
\cdot
\Bigl[\tfrac{1}{L'(\rho,\chi)\,w} - \tfrac{L''(\rho,\chi)}{2 L'(\rho,\chi)^2} + O(w)\Bigr].
\]

Collecting powers of \(w\):
- \(w^{-2}\) coefficient: \(\tfrac{1}{L'(\rho,\chi)}\) (from \(\tfrac{1}{w} \cdot \tfrac{1}{L'(\rho,\chi)\,w}\));
- \(w^{-1}\) coefficient: \(\tfrac{\log K}{L'(\rho,\chi)} - \tfrac{L''(\rho,\chi)}{2\,L'(\rho,\chi)^2}\) (cross-terms).

Hence
\begin{equation}
\label{eq:double-pole-residue}
\mathop{\mathrm{Res}}_{w = 0} F(w)
\;=\;
\frac{\log K}{L'(\rho,\chi)}
\;-\;
\frac{L''(\rho, \chi)}{2\,L'(\rho, \chi)^2}
\;=\;
\frac{\log K}{L'(\rho,\chi)} + C_1(\chi, \rho).
\end{equation}

This is \textbf{Lemma X.3.1 of the main text}, established algebraically;
the local part of (\ref{eq:cK-appendix}) is now done. It remains to
control the off-target zero residues, the trivial-zero residues, and
the contour pieces.

\subsection{Off-target zero residues, trivial-zero residues, and contour pieces}

We now shift the Perron contour from \(\mathrm{Re}(w) = \sigma_0' > \tfrac12\)
leftward to a zero-avoiding line \(\mathrm{Re}(w) = -A\) for some \(A > 0\),
crossing all the nontrivial zeros of \(L(s, \chi)\) in the strip
\(-A \le \mathrm{Re}(w) \le \sigma_0'\). The integrand \(F(w)\) has poles
at \(w = 0\) (the target double pole, with residue
(\ref{eq:double-pole-residue})), at \(w = \rho' - \rho\) for every
other nontrivial zero \(\rho' \ne \rho\) of \(L(s, \chi)\), at trivial
zeros, and at \(w = -\rho\) from a (potential) pole of \(L(s, \chi)\) at
\(s = 0\) --- but \(L(s, \chi)\) is entire for primitive non-principal
\(\chi\), so the pole at \(w = -\rho\) from \(1/L\) is absent. We do pick
up the term \(\mathop{\mathrm{Res}}_{s = 0}\bigl(x^s/(L(s,\chi)\,s)\bigr)\) from Inoue's formula
(this is the \(s = 0\) residue in the unshifted form; after dividing by
\(K^\rho\) it contributes a term of magnitude \(K^{-\rho}/(L(0,\chi)) =
O(K^{-1/2})\) in modulus, \(o(1)\)).

\subsubsection{Off-target nontrivial-zero residues}

For each \(\rho' \ne \rho\) a nontrivial zero of \(L(s, \chi)\), write
\(w_{\rho'} = \rho' - \rho\). Assuming \(\rho'\) is \emph{simple} (which is
generically expected and is the regime relevant for the leading +
subleading identity), the residue of \(F\) at \(w = w_{\rho'}\) is
\[
\mathop{\mathrm{Res}}_{w = w_{\rho'}} F(w)
\;=\;
\frac{K^{w_{\rho'}}}{w_{\rho'} \cdot L'(\rho', \chi)}
\;=\;
\frac{K^{\rho' - \rho}}{(\rho' - \rho)\,L'(\rho', \chi)}.
\]

Under RH for \(L(s, \chi)\), \(|K^{\rho' - \rho}| = 1\) for all
\(\rho' = \tfrac12 + i\gamma'\) with \(\gamma' \in \mathbb{R}\). The
aggregate off-target sum (over zeros \(\rho'\) in the strip
\(|\gamma'| \le T\)) is then bounded by
\[
\Bigl|\!\sum_{\rho' \ne \rho,\ |\gamma'| \le T}
\frac{K^{i(\gamma' - \tau)}}{(\rho' - \rho)\,L'(\rho', \chi)}\Bigr|
\;\ll\;
\sum_{\rho' \ne \rho,\ |\gamma'| \le T}
\frac{1}{|\gamma' - \tau|\,|L'(\rho', \chi)|}.
\]

This is the off-target \emph{aggregate} whose control determines the
\(o(1)\) rate.

\subsubsection{Trivial-zero residues}

Trivial zeros of \(L(s, \chi)\) are at \(s = -2k\) (\(\chi\) even) or
\(s = -2k - 1\) (\(\chi\) odd), \(k \ge 0\). After the shift \(w = s - \rho\),
these contribute residues of size \(K^{-\mathrm{Re}(\rho) - 2k}
= K^{-1/2 - 2k}\), summable to \(O(K^{-1/2})\), which is \(o(1)\).

\subsubsection{Shifted contour pieces}

The shifted contour \(\mathrm{Re}(w) = -A\) contributes, on a suitable
zero-avoiding truncation height \(T(K) \in [K/(\log K)^B, 2K/(\log K)^B]\), an
integral bounded by (using a reciprocal-\(L\) bound on the vertical
shifted line)
\[
\Bigl|\frac{1}{2\pi i}\int_{(-A)}\!\frac{K^w}{w\,L(w+\rho,\chi)}\,dw\Bigr|
\;\ll\;
K^{-A} \cdot T(K) \cdot \exp\bigl(C\,(\log\log T(K))^2\bigr).
\]

For \(A > 1/2\) and \(T(K) \le K/(\log K)^B\) with \(B\) sufficiently large,
this is \(K^{-A + 1 + o(1)} \cdot \exp(O((\log\log K)^2))\), which is
\(o(1)\).

Horizontal pieces give factors like \((T(K)/K)^A \cdot \mathrm{polylog}\),
\(o(1)\) for \(T(K) \le K (\log K)^{-B}\) and large enough \(B\).

\subsubsection{The Perron truncation error}

By Inoue~\cite[Theorem~1]{Inoue2021}, the truncation error
\(\mathcal{R}'(K, T(K), \rho)\) is bounded by
\[
\mathcal{R}'(K, T(K), \rho)
\;\ll\;
\frac{K^{\sigma_0' - 1/2} \log K}{T(K)}
\;+\;
\min\!\Bigl(1,\,\frac{K^{1/2}}{T(K) \langle K\rangle}\Bigr).
\]

For \(T(K) \asymp K (\log K)^{-B}\) and \(\sigma_0' = 1/2 + 1/\log K\), both
pieces are \(O(K^{-1/2 + \varepsilon}) \cdot (\log K)^{O(1)}\) for any
\(\varepsilon > 0\), in particular \(o(1)\).

\subsubsection{Assembly}

Combining the local residue (\ref{eq:double-pole-residue}) with the
off-target and contour pieces:
\begin{multline*}
c_K(\chi, \rho)
\;=\;
\frac{\log K}{L'(\rho, \chi)}
\;+\;
C_1(\chi, \rho)
\;+\;
\bigl(\text{off-target aggregate}\bigr)
\;+\;
O(K^{-1/2})
\\
\;+\;
O\!\bigl(K^{-A+1+o(1)}\exp(O((\log\log K)^2))\bigr)
\;+\;
\mathcal{R}'(K, T(K), \rho).
\end{multline*}
All terms in the last three parentheses are \(o(1)\). The off-target
aggregate is also \(o(1)\) under RH for \(L(s, \chi)\), by partial
summation against Soundararajan~\cite{Soundararajan2009}, \emph{Partial sums of the Möbius
function}, J. reine angew. Math. (Crelle's Journal) \textbf{631} (2009),
141--152, Theorem 1:
\[
\text{(\cite{Soundararajan2009})}
\quad
\text{under RH:}\quad
M(x) \;\ll\; \sqrt x \,\exp\!\bigl((\log x)^{1/2}\,(\log\log x)^{14}\bigr).
\]

Translating to \(M^*(x, \chi)\) (the character analogue) gives the
required \(o(1)\) rate for the off-target aggregate under RH for
\(L(s, \chi)\).

This establishes (\(\dagger\)). \(\hfill\square\)

\subsection{Rate of the \(o(1)\)}

The proof above yields the following rates.

{\def\LTcaptype{none} % do not increment counter
\begin{longtable}[]{@{}
  >{\raggedright\arraybackslash}p{(\linewidth - 2\tabcolsep) * \real{0.5000}}
  >{\raggedright\arraybackslash}p{(\linewidth - 2\tabcolsep) * \real{0.5000}}@{}}
\toprule\noalign{}
\begin{minipage}[b]{\linewidth}\raggedright
Rate
\end{minipage} & \begin{minipage}[b]{\linewidth}\raggedright
Hypotheses
\end{minipage} \\
\midrule\noalign{}
\endhead
\bottomrule\noalign{}
\endlastfoot
\(o(1) = O\bigl(K^{-1/2 + \varepsilon}\bigr)\) for every \(\varepsilon > 0\) & RH for \(L(s, \chi)\). \\
\(o(1) = O\bigl(K^{-1/2}\exp\bigl((\log K)^{1/2}(\log\log K)^{14}\bigr)\bigr)\) & RH for \(L(s, \chi)\), via the character analogue of Soundararajan~\cite[Theorem~1]{Soundararajan2009}. The operative rate for \S\,X.5: for \(\chi_{-4}, \chi_5, \chi_{11}\) the nontrivial zeros of \(L(s,\chi)\) in the relevant explicit-formula range are numerically verified on the critical line (Supplementary computation audit); not asserted as an unconditional theorem. \\
\(o(1) = O\bigl(\log K / \sqrt K\bigr)\) & RH plus a Gonek--Hejhal-type bound of the form \(\sum_{0 < \gamma < T} |L'(\rho, \chi)|^{-2} \ll T (\log T)^{O(1)}\). \\
\end{longtable}
}

In the numerical work of \S\,X.5.3, the empirical residual
\(R(K) := c_K - \log K / L'(\rho, \chi) - C_1(\chi, \rho)\) at
\(K = 2\cdot 10^6\) falls in the interval \([0.027,\,0.375]\) across the
four pairs; this is consistent with the Soundararajan-conditional rate
within an \(O(1)\) implicit-constant factor (specifically,
\(|R(K)| / (\log K / \sqrt K) \in [3,\,36]\), well within the typical
\(\exp\bigl(C\,(\log K)^{1/2}(\log\log K)^{14}\bigr)\) envelope at this
scale).

\subsection{What the proof does \emph{not} claim}

For clarity:

\begin{itemize}
\item
  The proof does \textbf{not} establish the Aoki--Koyama--Mertens limit
  for \(E_K \log K\), which is the \emph{multiplicative-side} asymptotic.
\item
  The proof does \textbf{not} establish the \emph{full} Perron-leading theorem
  \(c_K = \log K / L'(\rho) + o(\log K)\) unconditionally; the
  off-target aggregate's \(o(\log K)\) control (which is \emph{weaker} than
  the \(o(1)\) obtained in (\(\dagger\)) at the \emph{leading + subleading}
  level) depends on RH for \(L(s, \chi)\) and is what is required for
  composing with (AK) to obtain (NDC).
\item
  The proof does \textbf{not} address the possibility of off-target
  \emph{multiple} zeros, which under DRH/EDRH alone are not excluded;
  the manuscript's open Question Q:Perron asks for either
  the off-target multiple-zero residue control or a sufficient
  ``all crossed off-target zeros simple'' hypothesis.
\end{itemize}


\printbibliography

\end{document}
